\chapter{Troubleshooting}
\label{ch:trouble}

\section{The family or the counts look wrong}

The banner is the authority.
Compare it with Table~\ref{tab:counts}.

\begin{itemize}
\item \code{model_family=classical} together with a positive \code{n_qubits_value} means the run is classical: either you set \code{model_family = "classical"}, or a hybrid specification had zero qubits and was rewritten.
  The rewrite is printed by the stored family on the trial, which becomes \code{classical}.
\item \code{width=0} and \code{depth=0} on a full-quantum banner are expected.
  Those fields do not enter Equation~\eqref{eq:full-count}.
\item A hybrid total far above the classical total at the same width and depth is the encoder, the decoder, and the angles.
  Equation~\eqref{eq:hybrid-count} is the check.
\item Quantum parameters count only the circuit angles, \(3 m n_{q} n_{L}\) for \(m\) fields.
  Encoder and decoder weights are classical, including on a full-quantum model.
\end{itemize}

\section{Optuna rejects the study}

A message about a dynamic value space, a categorical distribution, or an incompatible distribution means the new suggestion does not match trials already stored under \code{study_name}.
The usual edits that cause it are adding or removing a family, adding or removing an embedding, or changing the min, max, or step of width, depth, qubits, or layers.

Point \code{storage_url} at a new file and set a new \code{study_name}.
Do not reuse the old name inside the new file either, if you copy the database by hand.
Confirm the path the trainer will open: a relative URL \code{sqlite:///NAME.db} is \file{analysis/NAME/NAME.db}.

\section{The search finishes without training the model you expected}

If \code{n_trials} completed trials are already in the study, a new process adds nothing and trains the current best trial.
Raise the cap, or use a new study.

If \code{use_optuna} is false, the tune list is empty, or \code{n_trials} is 0, the log says the search was skipped.
The model that trains is the fixed \code{*_value} architecture.
That is the behavior of Chapter~\ref{ch:examples}.

\code{--search-only} never trains the final model.
Run \code{--finalize} after the workers return.
\code{--finalize} and \code{--search-only} together are an error, and \code{--search-only} without a worker id is an error.

\section{Every quantum trial is discarded}

The parameter window is in force, and the sampled circuits do not fit.
A full-quantum model with many qubits and layers exceeds \code{total_params_max}.
A classical model smaller than \code{total_params_min} is discarded as well.
Read the discard line: it prints the count, the family, and the bounds.
Either widen the window, using Table~\ref{tab:counts} to recompute the overlap, or shrink \code{n_qubits_max} and \code{q_layers_max}.
Setting both bounds to \code{None} disables the window.
Do that in a new study if the old study's interpretation depended on the overlap.

\section{The loss columns}

Hard mode prints IC and BC as zero.
A nonzero IC or BC value in hard mode means the constraint string was not \code{hard}.
The legal strings are exactly \code{hard} and \code{soft}, in either letter case.

Soft mode can show a small PDE loss and a large penalty.
The Optuna objective includes the penalty.
A flat field can make a steady residual small while missing the boundary data; the penalty is what prevents that trial from winning.
If the penalty dominates for the whole trial, the boundary weights on the physical dataclass are large relative to the residual, or the network has not yet fit the trace.
The dam-break weights of 50 on the initial step and 10 on the walls are the reference balance for that problem.

Equal log-weights produce PDE weights \((1,1)\).
If the mass column is systematically smaller than the momentum column, that is the residual, not a hidden extra factor, as long as both log-weights are zero.

\section{Adam prints nothing}

\code{show_epoch_output} or \code{show_trial_epoch_output} is false, or the stride is larger than the epoch budget.
The final loss is still stored in the loss archive.
Trial rows are a separate print and still appear when a trial completes.

\section{Workers fail together}

\begin{itemize}
\item \code{One or more workers failed} from \file{run_cores.sh} means at least one \file{logs/worker_*.txt} ends in a traceback.
  Final training was not started.
\item \code{database is locked} that persists past the 120-second wait usually means a crashed process left a transaction, or two launches are sharing the file.
  Stop the extra launch.
  If no process has the database open, remove a stale \file{-wal} only when you have confirmed that no writer is alive.
  The trainer relocates a legacy database from the repository root or from \file{analysis/<problem name>/} into \file{analysis/<database name>/} on open.
  Two copies of the same file in those old locations are not merged; an existing destination is left in place.
\item A worker that cannot import PennyLane fails on its first quantum trial.
  A classical trial in the same study still completes.
  Install PennyLane and rerun; completed trials are kept.
\item An out-of-memory failure on a quantum trial is a large collocation grid times a large circuit.
  The circuit is evaluated as a statevector, batched over the collocation points.
  Reduce \code{trial_nx} and \code{trial_nt}, or reduce \code{n_qubits_max}, before raising the worker count.
\end{itemize}

\section{The analytic comparison retrains}

The script prints that it is training the best trial when the weight file is missing or does not match the architecture.
Matching requires the same family, sizes, learning rate, weights, and embedding.
A smoke test that overwrote \file{best_optuna_model.weights.h5} in a production folder is the usual cause.
Restore the weights from your own copy, or let the script train, which replaces the file with the architecture it selected.
Keep smoke tests on a separate \code{storage_url}, as in Chapter~\ref{ch:examples}.

The comparison selects the lowest \code{raw_objective}.
Final training selects Optuna's best objective.
Section on which trial is best, in Chapter~\ref{ch:search}, is the distinction.
Quote the trial number from the file you are reading.

\section{Full quantum is slow, and warmup did nothing}

Compilation happens on the first forward pass of each distinct circuit size in the process.
A search that draws many \((n_{q}, n_{L})\) pairs compiles many circuits.
That cost is paid again in every worker.

\code{hybrid_classical_warmup_fraction} does not apply to \code{full_quantum}.
The circuit is the model.
To keep a hybrid circuit off for the first quarter of a trial, set the fraction to \code{0.25} and confirm the warmup line in the trial log.
The score is always computed with the circuit on.

\section{Checklist before a long study}

\begin{enumerate}
\item The banner of a one-trial or fixed run shows the family, the embedding, and the parameter counts you computed.
\item \code{problem_name}, \code{constraints}, and the physical dataclass are the benchmark you mean to rank.
\item \code{storage_url} and \code{study_name} are new, or you intend to append to that study.
\item \code{model_family_options} lists only the families that should appear.
  A two-family preset is a preset that already removed \code{full_quantum}.
\item The parameter window matches the number of fields.
  The numbers 42 and 603 belong to the three-field eddy box.
\item \code{optuna_n_startup_trials} is above \code{n_trials} when you want equal family draws, and below it when you want TPE to take over.
\item \code{trial_epochs} is the budget you are willing to rank, and \code{final_adam_epochs} is the budget for the model you will plot.
\item \code{use_lbfgs} is the choice you want recorded in the summary.
\item PennyLane imports in a one-line hybrid or full-quantum build before you occupy a queue with hundreds of trials.
\end{enumerate}
